\section{Relaciones Principales del Sistema}

Las interacciones entre los diferentes elementos del sistema definen su comportamiento y flujo de trabajo. Estas relaciones se pueden clasificar en las siguientes categorías:

\subsection{Relaciones Jerárquicas y de Asignación}

\begin{itemize}
    \item Un \textbf{Supervisor de Mantenimiento} gestiona a uno o varios \textbf{Técnicos}.
    \item Un \textbf{Técnico de Mantenimiento} (o un \textbf{Contratista Externo}) es asignado a una o muchas \textbf{Órdenes de Trabajo}.
\end{itemize}

\subsection{Relaciones de Generación de Procesos (Causa-Efecto)}

\begin{itemize}
    \item Un \textbf{Reporte de Anomalía} (originado por una inspección) genera una \textbf{Orden de Trabajo} de tipo \textit{Correctivo}.
    \item Una \textbf{Alerta Predictiva} (originada por el modelo de IA) genera una \textbf{Orden de Trabajo} de tipo \textit{Predictivo}.
    \item Un \textbf{Plan de Mantenimiento Programado} genera \textbf{Órdenes de Trabajo} de tipo \textit{Preventivo}.
\end{itemize}

\subsection{Relaciones Funcionales y de Datos}

\begin{itemize}
    \item Un \textbf{Activo} posee un único \textbf{Historial de Mantenimiento}, el cual se compone de múltiples \textbf{Órdenes de Trabajo} cerradas.
    \item Un \textbf{Activo} está asociado a uno o varios flujos de \textbf{Datos de Sensores}.
    \item Una \textbf{Orden de Trabajo} consume uno o varios \textbf{Repuestos} del \textbf{Inventario (MRO)}.
    \item El \textbf{Gerente de Planta} consulta los \textbf{Informes de KPIs}, los cuales son calculados a partir de los datos agregados de todas las \textbf{Órdenes de Trabajo} completadas.
\end{itemize}
